% SGA 3, Expose V, Loop 1 English reconstruction.
% Staging scope only: component-local PDF pp. 32--35 inclusive.
% Authority: Polo--Gille Exp5-13oct24.pdf, SHA-256
% 9198200633F929FE1822520371A9200DA4F8CF2513EFAB3D6A0C7E9330DB84CF.
% Locator mapping:
%   local p. 32 = combined-reader p. 294 = re-edited header p. 282.
%   local p. 33 = combined-reader p. 295 = re-edited header p. 283;
%     original-edition marginal marker 286 occurs after the algebra diagram.
%   local p. 34 = combined-reader p. 296 = re-edited header p. 284.
%   local p. 35 = combined-reader p. 297 = re-edited header p. 285.
% Jacob Reinhold's CC BY 4.0 Markdown is comparison/drafting material only.

\phantomsection\label{thm:sga3-V-9.0}%
\noindent\textbf{定理 9.0.}---%
\textit{设 \(S\) 为任意概形, 并设}
\begin{equation*}
  \resizebox{0.34\textwidth}{!}{%
    \SGAThreeNativeInput{figures/exp5/exp5_localp032_theorem90_groupoid.png}}
  \phantomsection\label{fig:sga3-V-9.0-groupoid}
\end{equation*}
\textit{为一个 \((\mathbf{Sch}/S)\)-群胚, 使得 \(X_0\) 在 \(S\) 上
有限呈示且拟射影, \(d_1\) 有限
呈示、固有且平坦, 并且态射}
\[
  d_0\boxtimes d_1:X_1\longrightarrow X_0\times_S X_0
\]
\textit{拟有限. 则:}
\begin{enumerate}
\item[(1)]\phantomsection\label{thm:sga3-V-9.0-1}%
\textit{\(X_0\) 的每一点 \(x\) 都有一个饱和开
邻域 \(W\), 使得 \(X_\ast\) 在
\(W\) 上诱导的群胚具有一个拟截面.}

\item[(2)]\phantomsection\label{thm:sga3-V-9.0-2}%
\textit{设 \((Y,p)\) 为所有赋环空间的范畴中
\((d_0,d_1)\) 的余核. 则 \(Y\) 是概形, \(p\) 是
概形态射, 且 \((Y,p)\) 是
\((\mathbf{Sch}/S)\) 中 \((d_0,d_1)\) 的余核.}

\item[(3)]\phantomsection\label{thm:sga3-V-9.0-3}%
\textit{态射 \(p\) 满射、为开映射且普遍
闭.}

\item[(4)]\phantomsection\label{thm:sga3-V-9.0-4}%
\textit{以 \(d_0\) 和 \(d_1\) 为分量的态射}
\[
  d:X_1\longrightarrow X_0\times_Y X_0
\]
\textit{满射.}

\item[(5)]\phantomsection\label{thm:sga3-V-9.0-5}%
\textit{若 \((d_0,d_1)\) 是等价对, 则:}
  \begin{enumerate}
  \item[(a)]\phantomsection\label{thm:sga3-V-9.0-5-a}%
  \textit{态射}
  \[
    d:X_1\longrightarrow X_0\times_Y X_0
  \]
  \textit{是同构, 且 \(p\) 忠实平坦;}

  \item[(b)]\phantomsection\label{thm:sga3-V-9.0-5-b}%
  \textit{态射 \(p\) 和 \(Y\to S\) 有限
  呈示, 且 \((Y,p)\) 是 fppf 拓扑下的层范畴中
  \((d_0,d_1)\) 的余核.}
  \end{enumerate}
\end{enumerate}

\noindent\textit{证明.} 对于 \hyperref[thm:sga3-V-9.0-1]{(1)}, 此问题在 \(S\) 上是局部的,
故可以假设 \(S=\operatorname{Spec}B\) 仿射. 此时
存在一个在 \(\mathbb Z\) 上有限型的环 \(A\), 一个态射
\[
  S\longrightarrow T=\operatorname{Spec}A,
\]
以及一个 \((\mathbf{Sch}/T)\)-群胚 \(Z_\ast\), 使得
\[
  X_\ast\simeq Z_\ast\times_T S.
\]
事实上, 对
\(S_0=\operatorname{Spec}\mathbb Z\) 和
\(S_i=\operatorname{Spec}A_i\) 应用 EGA~IV\textsubscript{3}, 8.8.3,
其中 \(A_i\) 遍历 \(B\) 的有限生成
\(\mathbb Z\)-子代数. 此外, 由
EGA~IV\textsubscript{3}, 8.10.5, 可以假设 \(Z_\ast\)
满足定理~\hyperref[thm:sga3-V-7.1]{7.1} 的假设.
因此 \(Z_\ast\) 局部具有
拟截面.

由 \hyperref[par:sga3-V-9-a]{(a)} 部分, \(X_\ast\) 也有同样的性质,
而断言 \hyperref[thm:sga3-V-9.0-2]{(2)}, \hyperref[thm:sga3-V-9.0-3]{(3)}, \hyperref[thm:sga3-V-9.0-4]{(4)} 和 \hyperref[thm:sga3-V-9.0-5-a]{(5)(a)}
由引理~\hyperref[lem:sga3-V-6.1]{6.1} 得到,
与定理~\hyperref[thm:sga3-V-7.1]{7.1} 的证明相同.

\medskip
\phantomsection\label{proof:sga3-V-9.0-c}%
\noindent\textbf{c) \hyperref[thm:sga3-V-9.0-5-b]{(5)(b)} 的证明.}---%
下面证明 \(Y\to S\) 有限
呈示.\footnote{\label{note:sga3-V-47}编者注: 原文说这由
下面的命题~\hyperref[sga3:v:proposition:9:1]{9.1} 得到.
我们未能重构该论证.
下面的证明由 O.~Gabber 提供.}
由假设, \((d_0^{X_\ast},d_1^{X_\ast})\) 是等价
对; 即
\[
  d^{X_\ast}:X_1\longrightarrow X_0\times_S X_0
\]
是单态射. 由 EGA~IV\textsubscript{3}, 8.10.5, 必要时扩大
\(A\) 后, 可以假设
\[
  d^{Z_\ast}:Z_1\longrightarrow Z_0\times_T Z_0
\]
是单态射. 由于 \(T=\operatorname{Spec}A\), 其中 \(A\)
为 Noether 环, 定理~\hyperref[thm:sga3-V-7.1]{7.1} 表明,
群胚 \(Z_\ast\) 在
\((\mathbf{Sch}/T)\) 中有余核 \((Q,q)\), \(q\) 和 \(Q\to T\) 有限
呈示, 而且
\[
  q:Z_0\longrightarrow Q
\]
忠实平坦, \(d^{Z_\ast}\) 诱导同构
\[
  Z_1\xrightarrow{\ \sim\ }Z_0\times_Q Z_0.
\]
令 \(Q_S=Q\times_T S\). 由于 \(X_i\simeq Z_i\times_T S\), 得到
同构
\[
  d^{Z_\ast}\times_T S:
  X_1\xrightarrow{\ \sim\ }X_0\times_{Q_S}X_0.
\]
用 \(\phi\) 表示其逆, 并设 \(\pi\) 为典范
态射
\[
  X_0\times_Y X_0\longrightarrow X_0\times_{Q_S}X_0.
\]

则 \(\phi\circ\pi\) 是
\[
  d_0\boxtimes d_1:
  X_1\xrightarrow{\ \sim\ }X_0\times_Y X_0.
\]
的逆. 因此 \(X_\ast\) 定义的等价关系, 即
图~\hyperref[fig:sga3-V-9.0-equivalence-relation]{\ref*{fig:sga3-V-9.0-equivalence-relation}}
所示的单态射,
\SGAThreeDiagram[0.78\textwidth]
  {figures/exp5/exp5_localp033_equivalence_relation.png}
  {\(X_\ast\) 定义的等价关系单态射.}
  {fig:sga3-V-9.0-equivalence-relation}
  {第 V 篇, 组件局部第~33 页 / 合订阅读版第~295 页 / 修订版页眉第~283 页}
可与 \(q_S:X_0\to Q_S\) 定义的等价关系
\(\mathscr R(q_S)\) 等同. 由于 \(q_S\) 忠实平坦且有限
呈示, 因而是普遍有效满态射,
\(\mathscr R(q_S)\) 的商为 \(Q_S\); 见 \SGARefLink{sga3:iv:definition:3:3:2}{第 IV 篇, 3.3.2}.
因此,
\[
  Y\simeq Q\times_T S,
\]
故 \(Y\to S\) 和 \(p\) 有限呈示. 此外, 由
\SGARefLink{sga3:iv:corollary:6:3:3}{第 IV 篇, 6.3.3}, \((Y,p)\) 也是 fppf 拓扑下的层范畴中
\((d_0,d_1)\) 的余核.

\medskip
\SGARefTarget{sga3:v:proposition:9:1}
\phantomsection\label{prop:sga3-V-9.1}%
\noindent\textbf{命题 9.1.}---%
\textit{考虑概形态射}
\[
  X_0\xrightarrow{\,p\,}Y\xrightarrow{\,q\,}S
\]
\textit{使得 \(qp\) 有限型 (相应地, 有限
呈示), 且 \(p\) 忠实平坦且有限
呈示. 则 \(q\) 有限型 (相应地, 有限
呈示).}%
\begingroup
\renewcommand{\thefootnote}{*}%
\footnote{\label{note:sga3-V-9.1-generalization}更一般的结果见
EGA~IV\textsubscript{4}, 17.7.5.}
\endgroup
\addtocounter{footnote}{-1}%

\noindent 由于 \(p\) 满射且 \(qp\) 拟紧,
\(q\) 拟紧. 因此可以假设 \(S\), \(Y\)
和 \(X_0\) 都仿射, 相应的环分别为 \(A\), \(B\) 和 \(C\).
有
\[
  B=\varinjlim_i B_i,
\]
其中 \(B_i\) 遍历 \(B\) 的有限生成
\(A\)-子代数. 由于 \(C\) 在 \(B\) 上有限呈示,
存在一个指标 \(i_0\)、一个有限呈示的 \(B_{i_0}\)-代数
\(C_{i_0}\) 以及一个同构
\[
  C\simeq C_{i_0}\otimes_{B_{i_0}}B.
\]
对 \(i\geq i_0\), 令
\[
  C_i=C_{i_0}\otimes_{B_{i_0}}B_i;
\]
则
\[
  C\simeq C_i\otimes_{B_i}B.
\]

\SGAThreeDiagram[0.22\textwidth]
  {figures/exp5/exp5_localp033_algebra_limit_diagram.png}
  {极限论证中使用的相容 \(A\)-代数.}
  {fig:sga3-V-9.1-algebra-limit}
  {第 V 篇, 组件局部第~33 页 / 合订阅读版第~295 页 / 修订版页眉第~283 页}

由于 \(C\) 在 \(B\) 上忠实平坦, EGA~IV\textsubscript{3},
11.2.6 和 8.10.5(vi) 给出一个指标 \(i_1\geq i_0\), 使得
\(C_{i_1}\) 在 \(B_{i_1}\) 上忠实平坦. 因此,
对每个 \(i\geq i_1\), \(C_i\) 在 \(B_i\) 上忠实平坦.
对 \(i\geq i_1\), 典范映射 \(C_i\to C\) 因而是单射,
因为它由 \(B_i\to B\) 经忠实平坦基
扩张得到.

% The re-edited source places original-edition marginal marker 286 here.

若 \(C\) 在 \(A\) 上有限型, 则存在
一个指标 \(j\geq i_1\), 使得 \(C_j=C\), 从而 \(B_j=B\), 因为
\(C_j\) 在 \(B_j\) 上忠实平坦. 因此 \(B\) 在 \(A\) 上有限型.

现在假设 \(C\) 在 \(A\) 上有限呈示. 由上述结果,
\(B\) 在 \(A\) 上有限型, 因而具有形式
\[
  B=\overline B/I,
\]
其中 \(\overline B\) 是 \(A\) 上有限个未定元的多项式代数,
\(I\) 是 \(\overline B\) 的理想.
理想 \(I\) 是其有限生成子理想
\(I_\alpha\) 的并; 因此,
\[
  B=\varinjlim_\alpha B_\alpha,
  \qquad B_\alpha=\overline B/I_\alpha.
\]
如上论证, 存在一个指标 \(\alpha_0\)、一个
有限呈示的 \(B_{\alpha_0}\)-代数 \(C_{\alpha_0}\),
以及一个同构
\[
  C\simeq C_{\alpha_0}\otimes_{B_{\alpha_0}}B.
\]
对 \(\alpha\geq\alpha_0\), 令
\[
  C_\alpha=C_{\alpha_0}\otimes_{B_{\alpha_0}}B_\alpha,
\]
于是
\[
  C\simeq C_\alpha\otimes_{B_\alpha}B
  \qquad(\alpha\geq\alpha_0).
\]
再由 EGA~IV\textsubscript{3}, 11.2.6 及 8.10.5(vi),
对所有充分大的 \(\alpha\), \(C_\alpha\) 在 \(B_\alpha\) 上忠实
平坦. 此时, 映射
\(C_\alpha\to C\) 和 \(C_\alpha\to C_\beta\) 的核分别可与
\[
  C_\alpha\otimes_{B_\alpha}(I/I_\alpha)
  \quad\text{和}\quad
  C_\alpha\otimes_{B_\alpha}(I_\beta/I_\alpha)
  \qquad(\beta\geq\alpha).
\]
等同.

由于 \(C_\alpha\) 和 \(C\) 在 \(A\) 上有限呈示,
且 \(C_\alpha\to C\) 满射,
\[
  C_\alpha\otimes_{B_\alpha}(I/I_\alpha)
\]
是有限生成理想,\footnote{\label{note:sga3-V-48}编者注:
见 EGA~IV\textsubscript{1}, 1.4.4.} 并且它是下列理想的并:
\[
  C_\alpha\otimes_{B_\alpha}(I_\beta/I_\alpha).
\]
因此, 对所有充分大的 \(\beta\),
\[
  C_\alpha\otimes_{B_\alpha}(I_\beta/I_\alpha)
  =
  C_\alpha\otimes_{B_\alpha}(I/I_\alpha).
\]
由于 \(C_\alpha\) 在 \(B_\alpha\) 上忠实平坦, 这还给出
\(I_\beta=I\). 因此 \(B\) 在 \(A\) 上有限呈示.

\bigskip
\phantomsection\label{sec:sga3-V-10}%
\section*{10. 补充: 模群概形的商}
\addcontentsline{toc}{section}{10. 补充: 模群概形的商}

下面的 \hyperref[sec:sga3-V-10.2]{10.2}--\hyperref[sec:sga3-V-10.4]{10.4} 节根据
M.~Raynaud 的建议写成, 将上述定理应用于
群概形作用的情形. 为便于读者查阅, \hyperref[sga3:v:paragraph:10-1]{10.1 节}先
重述第 XVI 篇的陈述 2.1--2.3.

\medskip
\phantomsection\label{subsec:sga3-V-10.1}%
\noindent\SGAInlineRefTarget{sga3:v:paragraph:10-1}\textbf{10.1. 商的可表性定理.}---%
先回顾\hyperref[thm:sga3-V-10.1.1]{下述结果}.

\medskip
\phantomsection\label{thm:sga3-V-10.1.1}%
\noindent\textbf{定理 10.1.1.}---%
\textit{设 \(S\) 为概形, \(X\) 和 \(Y\) 为两个
\(S\)-概形, \(f:X\to Y\) 为 \(S\)-态射. 假设
下列两个条件之一成立:}
\begin{enumerate}
\item[\(\alpha\))]\phantomsection\label{thm:sga3-V-10.1.1-alpha}%
\textit{态射 \(f\) 局部有限呈示;}
\item[\(\beta\))]\phantomsection\label{thm:sga3-V-10.1.1-beta}%
\textit{概形 \(S\) 局部 Noether, 且 \(X\) 在 \(S\) 上局部
有限型.}
\end{enumerate}
\textit{则下列条件等价:}
\begin{enumerate}
\item[(i)]\phantomsection\label{thm:sga3-V-10.1.1-i}%
\textit{存在一个 \(S\)-概形 \(X'\) 及一个分解}
\[
  f:X\xrightarrow{\,f'\,}X'\xrightarrow{\,f''\,}Y,
\]
\textit{其中 \(f'\) 为忠实平坦且局部有限呈示的
\(S\)-态射, \(f''\) 为单态射;}

\item[(ii)]\phantomsection\label{thm:sga3-V-10.1.1-ii}%
\textit{第一投影}
\[
  p_1:X\times_Y X\longrightarrow X
\]
\textit{平坦.}
\end{enumerate}
\textit{此外, 若这些条件成立, 则 \((X',f')\) 为 fppf 拓扑下
\(X\) 模 \(f\) 所定义的等价关系的商.
因此 \hyperref[thm:sga3-V-10.1.1-i]{(i)} 中的分解 \(f=f''\circ f'\)
在同构意义下唯一.}

\(Y\) 局部 Noether 且 \(X\) 在 \(Y\) 上有限
型的情形已在 \hyperref[bib:sga3-V-Mur65]{[Mur65]} 的定理~2 的推论~2 中处理.
我们将看到一般情形可化归于该情形.

先作几点说明.

\medskip
\phantomsection\label{par:sga3-V-10.1.1-a}%
\noindent\textbf{a)} 蕴涵 \hyperref[thm:sga3-V-10.1.1-i]{(i)}\(\Rightarrow\)\hyperref[thm:sga3-V-10.1.1-ii]{(ii)}
显然成立. 事实上, 第一投影
\[
  p'_1:X\times_{X'}X\longrightarrow X
\]
经由 \(X\times_Y X\) 分解:
\[
  p'_1:
  X\times_{X'}X\xrightarrow{\,u\,}X\times_Y X
  \xrightarrow{\,p_1\,}X.
\]
态射 \(u\) 为同构, 因为 \(f''\) 为
单态射; 而 \(p'_1\) 平坦, 因为 \(f'\) 平坦;
故 \(p_1\) 平坦.

\medskip
\phantomsection\label{par:sga3-V-10.1.1-b}%
\noindent\textbf{b)} 定理~\hyperref[thm:sga3-V-10.1.1]{10.1.1} 的断言
在 \(Y\) 上是局部的, 因而在 \(S\) 上也是局部的.
它们在 \(X\) 上也是局部的, 这很容易由局部有限呈示的平坦态射
为开映射这一事实得到
(EGA~IV\textsubscript{3}, 11.3.1).

\medskip
\phantomsection\label{par:sga3-V-10.1.1-c}%
\noindent\textbf{c)} 在定理~\hyperref[thm:sga3-V-10.1.1]{10.1.1} 的假设
\(\hyperref[thm:sga3-V-10.1.1-alpha]{\alpha}\) 下,
上述说明将我们化归于 \(X\) 和 \(Y\) 仿射且 \(f\)
有限呈示的情形. 必要时以 \(Y\) 代替 \(S\),
可以假设 \(X\) 和 \(Y\) 在 \(S\) 上有限呈示.
于是 EGA~IV\textsubscript{3}, 11.2.6 将我们化归于
\(S\) 为 Noether 概形的情形.

\medskip
\phantomsection\label{par:sga3-V-10.1.1-d}%
\noindent\textbf{d)} 在定理~\hyperref[thm:sga3-V-10.1.1]{10.1.1} 的假设
\(\hyperref[thm:sga3-V-10.1.1-beta]{\beta}\) 下,
可以假设 \(S\), \(X\) 和 \(Y\) 仿射, \(S\) 为 Noether 概形,
且 \(X\) 在 \(S\) 上有限型. 将 \(Y\) 写成
\(S\) 上有限型仿射概形 \(Y_i\) 的滤过
逆极限. 概形
\[
  X\times_{Y_i}X
\]
构成 \(X\times_S X\) 的一个滤过递减闭子概形族,
其逆极限为 \(X\times_Y X\). 由于
\(X\times_S X\) 为 Noether 概形, 对所有充分大的 \(i\),
有
\[
  X\times_{Y_i}X=X\times_Y X.
\]
因此,
\[
  f_i:X\xrightarrow{\,f\,}Y\longrightarrow Y_i
\]
在 \(f\) 满足定理~\hyperref[thm:sga3-V-10.1.1]{10.1.1} 的
\hyperref[thm:sga3-V-10.1.1-ii]{条件 (ii)} 时也满足该条件.
\(f\) 在 \(X\) 上定义的等价关系与
\(f_i\) 定义的等价关系相同. 因此只需对 \(f_i\) 证明
\hyperref[thm:sga3-V-10.1.1-ii]{(ii)}\(\Rightarrow\)\hyperref[thm:sga3-V-10.1.1-i]{(i)}, 从而化归于
\(Y\) 在 \(S\) 上有限型的情形.

\medskip
\noindent\textit{对群概形的应用.}
设 \(S\) 为概形, \(G\) 为在 \(S\) 上局部有限
呈示的 \(S\)-群概形, 并假设 \(G\) 左作用于
\(S\)-概形 \(X\). 若 \(X\to S\) 有一个截面 \(\xi\),
回顾稳定化子 \(\operatorname{Stab}_G(\xi)\) 可由
\(G\) 的子群概形表出; 见 \SGARefLink{sga3:i:definition:2:3:3}{第 I 篇, 2.3.3}.

\medskip
\phantomsection\label{thm:sga3-V-10.1.2}%
\noindent\textbf{定理 10.1.2.}---%
\textit{设 \(S\) 为概形, \(G\) 为在 \(S\) 上局部有限呈示的
\(S\)-群概形, 作用于 \(S\)-概形
\(X\). 假设 \(X\to S\) 有一个截面 \(\xi\), 其在 \(G\) 中的稳定化子
\(H\) 在 \(S\) 上平坦. 若下列假设之一成立:}
\begin{enumerate}
\item[\textit{a)}]\phantomsection\label{thm:sga3-V-10.1.2-a}%
\textit{\(X\) 在 \(S\) 上局部有限型;}
\item[\textit{b)}]\phantomsection\label{thm:sga3-V-10.1.2-b}%
\textit{\(S\) 局部 Noether,}
\end{enumerate}
\textit{则 fppf 商层 \(G/H\) 可由一个
\(S\)-概形表出, 它在 \(S\) 上局部有限呈示, 且
\(S\)-态射}
\[
  f:G\longrightarrow X,\qquad g\longmapsto g\cdot\xi
\]

% Exact continuation cursor: component-local p. 36 / combined-reader p. 298
% / re-edited header p. 286, beginning "factors as".
% SGA3-V-V2-EOL-PRESERVATION

% 翻译统计: 源行数 437; 输出行数 433 (含本统计行); 行内/行间数学片段 268; 翻译占位 0. 底本疑点另报主线程.
